% Checkpoint38 isolated proposal. All checkpoint37 bytes remain unchanged.
% Complete source-gauge proof from A's 1d82 note and DE's ae136 fixed-group proof.
\statement{A.135}{Proved}{The literal annular source gauge and its trivial fixed subgroup}
Let \(\alpha=(-1,\phi)\) be the ORIGINAL negative full-GT involution,
with its canonical commutator norm word \(u\). Retain the actual
completed annular source and its two based maps from A.134:
\[
 \begin{gathered}
 T=\widehat{(B_4/Z(B_4))\times_{S_4}S_3},\\
 \jmath:T\longrightarrow G,\qquad q:T\longrightarrow B.
 \end{gathered}
 \tag{A.135.1}
\]
Here \(P_T=\ker(T\to S_3)\), \(B=C_2(s)\amalg C_3(r)\), and
\[
 \begin{array}{c|ccc}
       &t&\sigma&\tau\\\hline
 \jmath&X&\sigma&\tau\\
 q&1&r^{-1}s&sr^{-1}.
 \end{array}
 \tag{A.135.2}
\]
Put \(t_2=\sigma t\sigma\), \(K_S=\sigma\tau^2\sigma\), and
\[
 \begin{gathered}
 \Phi_S=\phi(t,K_S),\qquad Z_\phi=\phi(t_2,\tau^2),\qquad
 h_S=\Phi_S Z_\phi,\\
 U_S=u(t,K_S),\qquad Z_u=u(t_2,\tau^2),\qquad r_S=U_SZ_u.
 \end{gathered}
 \tag{A.135.3}
\]
The actual source involution \(\theta_T\) satisfies the literal
source identities
\[
 \begin{gathered}
 \theta_T(t_2)=t_2^{-1},\qquad
 \theta_T(K_S)=h_S^{-1}K_S^{-1}h_S,\\
 \theta_T(r_S)=h_S^{-1}r_S,\qquad
 h_S=r_S\theta_T(r_S)^{-1}.
 \end{gathered}
 \tag{A.135.4}
\]
Consequently the ORIGINAL source gauge is the genuine involution
\[
 \beta_T=\operatorname{Ad}(h_S)\theta_T
 =\operatorname{Ad}(r_S)\theta_T\operatorname{Ad}(r_S^{-1}).
 \tag{A.135.5}
\]
Its domains and both image actions are exactly
\[
 \begin{gathered}
 \jmath\theta_T=\Theta_0\jmath,
 \quad q\theta_T=\alpha q,
 \quad \Theta_0=\operatorname{Ad}(h_1^{-1})\Theta,\\
 \jmath\beta_T=\Theta\jmath,
 \qquad q\beta_T=\operatorname{Ad}(\phi)\alpha q,
 \qquad \jmath(h_S)=h_1,\quad q(h_S)=\phi.
 \end{gathered}
 \tag{A.135.6}
\]
All four full fixed subgroups are trivial:
\[
 P_T^{\theta_T}=T^{\theta_T}
       =P_T^{\beta_T}=T^{\beta_T}=1.
 \tag{A.135.7}
\]
In particular the actual closed annular image \(H=\jmath(T)\),
using A.134's proved full embedding, satisfies
\[
 H^{\Theta_0}=H^\Theta=1.
 \tag{A.135.8}
\]
Neither of the two specified nontrivial Y-direction cubics belongs
to this image:
\[
 \begin{gathered}
 \Pi_0\notin H,\qquad
 \Pi_u^*=B_u\Pi_u B_u^{-1}=r_u\Pi_0r_u^{-1}\notin H,\\
 r_u=\jmath(r_S),\qquad B_u=u(\sigma^2,\tau^2).
 \end{gathered}
 \tag{A.135.9}
\]
These assertions neither make the specified positive source cube
fixed nor prove its vanishing, the original boundary equation,
or the original CM character value.

\begin{proof}
\emph{1. The actual source action and the original pentagon.}\par
Before completion, the source has the type-B presentation
\[
 \begin{gathered}
 \sigma\tau\sigma=\tau\sigma\tau,
 \qquad [t,\tau]=1,
 \qquad (t\sigma)^2=(\sigma t)^2,\\
 R=t\sigma\tau,\qquad R^3=1.
 \end{gathered}
 \tag{A.135.10}
\]
Indeed \(R\) is conjugate to \(\tau\sigma t\), whose cube is the
central full twist before the central quotient. The source and map
are those in A.121 and \citep[Proposition~2.20(1)]{A_Enriquez2014}.
The full genus-zero action is the existing A.36--A.37/C.28 action;
its original normalized rows restrict to
\[
 \theta_T(t)=t^{-1},\quad
 \theta_T(\sigma)=A^{-1}\sigma^{-1}A,\quad
 \theta_T(\tau)=Z_\phi^{-1}\tau^{-1}Z_\phi,
 \qquad A=\phi(t,\sigma^2).
 \tag{A.135.11}
\]
It preserves the permutation quotient and inverts the central full
twist. It is therefore an actual involution of the indicated source,
not a substitution inferred from its elliptic image.

Write
\[
 \begin{gathered}
 a=\sigma^2,\quad b=\tau^2,\quad
 t_3=\tau t_2\tau,\quad K_{2,S}=\tau\sigma^2\tau,\\
 F=\phi(a,b),\qquad E=\phi(t_3^{-1},K_{2,S}).
 \end{gathered}
\]
The original four-strand pentagon, evaluated in the actual type-B
source, is the complete identity
\[
                         FEA=\Phi_S Z_\phi=h_S.
 \tag{A.135.12}
\]
This is A.121.5's marked type-B pentagon. Any central decoration
before the central quotient disappears from a commutator word.
The parameter here is the ORIGINAL \(\phi\), not a prospective
positive GT parameter.

The type-B relations give \([t,t_2]=[t,b]=1\). The braid relation
and \([t,\tau]=1\) give \([\sigma,t_3]=1\). In the three-strand
braid subgroup the word
\[
 Z_c=(\sigma\tau)^3=aK_{2,S}=bK_S
 \tag{A.135.13}
\]
is central. These exact facts imply
\[
 [t,Z_\phi]=1,\qquad [\sigma,E]=1,\qquad [b,\Phi_S]=1.
 \tag{A.135.14}
\]
No centrality of \(Z_c\) in the entire source is used.

\emph{2. The complete source row for \(K_S\).}\par
The actual normalized negative GT action on the abstract
three-strand braid group sends
\(\sigma\mapsto\sigma^{-1}\) and
\(\tau\mapsto F^{-1}\tau^{-1}F\). It inverts \(Z_c\).
Conjugating this action by \(F\) gives the action \(\nu\) with
\[
 \nu(\sigma)=F\sigma^{-1}F^{-1},\quad
 \nu(\tau)=\tau^{-1},\quad
 \nu(Z_c)=Z_c^{-1}.
\]
Since \(K_S=Z_cb^{-1}\), it follows that
\[
 F\sigma^{-1}F^{-1}b^{-1}F\sigma^{-1}F^{-1}
       =\nu(K_S)=Z_c^{-1}b=K_S^{-1}.
 \tag{A.135.15}
\]
This abstract braid identity may be evaluated in the source; an
embedding of that braid subgroup is not required.

Equation (A.135.12) gives
\(AZ_\phi^{-1}=E^{-1}F^{-1}\Phi_S\). Substitute the actual
rows (A.135.11), preserving their order:
\[
 \begin{aligned}
 h_S\theta_T(K_S)h_S^{-1}
 &=FE\sigma^{-1}(AZ_\phi^{-1})b^{-1}
       (Z_\phi A^{-1})\sigma^{-1}E^{-1}F^{-1}\\
 &=F\sigma^{-1}F^{-1}\Phi_S b^{-1}\Phi_S^{-1}
       F\sigma^{-1}F^{-1}\\
 &=F\sigma^{-1}F^{-1}b^{-1}F\sigma^{-1}F^{-1}
 =K_S^{-1}.
 \end{aligned}
 \tag{A.135.16}
\]
The middle cancellations are exactly (A.135.14); the last step is
(A.135.15). This proves the claimed row in the full source.

\emph{3. The other exact pair and both canonical norm evaluations.}\par
The same original \(\phi\) satisfies the negative three-cycle
and the multiplier-one three-cycle, the latter because the actual
composite \(\alpha c\) is \((1,\phi)\). Put
\(c_a=(ta)^{-1}\) and \(J_a=\phi(a,c_a)\). The zero cycle and
duality give \(\phi(c_a,t)=A^{-1}J_a^{-1}\); substituting in the
negative cycle gives
\[
                    At^{-1}A^{-1}=aJ_a^{-1}c_aJ_a.
 \tag{A.135.17}
\]
In the subgroup generated by \(t,\sigma\), the word
\(t t_2=(t\sigma)^2\) is central and
\(\sigma c_a\sigma^{-1}=t_2^{-1}\). Therefore removing this
commuting central decoration from the commutator word gives
\[
 \sigma J_a^{-1}\sigma^{-1}
   =\phi(a,t_2^{-1})^{-1}
   =\phi(t_2^{-1},a)=A.
\]
Multiplying (A.135.17) on both sides by \(\sigma^{-1}\) gives
\[
 \sigma^{-1}At^{-1}A^{-1}\sigma^{-1}
       =At_2^{-1}A^{-1}.
\]
Apply (A.135.11) to \(t_2=\sigma t\sigma\) to obtain
\[
 \theta_T(t_2)=t_2^{-1},\qquad
 \theta_T(b)=Z_\phi^{-1}b^{-1}Z_\phi.
 \tag{A.135.18}
\]
This is a source word calculation using the original cycles; it is
not a positive hexagon assumption on \(u\).

Since \([Z_\phi,t]=1\), the action on the pair \((t,K_S)\) is
the original normalized negative action conjugated by
\(Z_\phi^{-1}\). The actual canonical identity
\(\alpha(u)=\phi^{-1}u\) may thus be evaluated on the two
specified pairs, without any free-pair assumption:
\[
 \begin{aligned}
 \theta_T(U_S)
   &=Z_\phi^{-1}\Phi_S^{-1}U_SZ_\phi
     =h_S^{-1}U_SZ_\phi,\\
 \theta_T(Z_u)&=Z_\phi^{-1}Z_u.
 \end{aligned}
 \tag{A.135.19}
\]
Multiplication gives \(\theta_T(r_S)=h_S^{-1}r_S\), and hence
the exact identity \(r_S\theta_T(r_S)^{-1}=h_S\).
Every argument in these two evaluations is pure, so
\(r_S,h_S\in P_T'\). This proves (A.135.4)--(A.135.5),
including the fact that the gauge is a conjugate involution.

The given source maps send \(r_S\) to A.133's \(r_u\) and
\(h_S\) to A.132's \(h_1\). Under the original forgetting map,
\(q(\Phi_S)=1\), \(q(Z_\phi)=\phi\), and \(q(r_S)=u\).
Together with \(q\theta_T=\alpha q\) these facts give all of
(A.135.6). In the external involution extension, the ORIGINAL
source gauge consequently maps its involutory symbol to
\(\phi\alpha\) in the original modular extension, not to
\(\alpha\) without this correction.

\emph{4. The exact C.28 input and the central-quotient transfer.}\par
C.28 proves trivial geometric centralizers for every genuine pure
negative involution of the full profinite free group of a punctured
sphere, with each individual peripheral class inverted. Its
full-coefficient argument excludes pro-odd centralizers as well.
For the actual full-GT action, C.28 verifies these hypotheses in
the full pure braid tower: finite point-pushing cusps have multiplier
\(-1\), and the infinity cusp is the quotient of the corresponding
two full twists. The split Fadell--Neuwirth sequences and
characteristic finite-index free-kernel refinements give their full
profinite exactness. Induction from the inverted two-strand group
therefore gives
\[
                  (\widehat P_4)^{\widetilde\theta}=1.
 \tag{A.135.20}
\]
This is the actual full-GT action on the whole tower. No five-point
action of a merely gentle pair is presumed.

A quotient need not inherit a trivial fixed subgroup. Here there
is an explicit equivariant splitting. The central full twist
\(z_4\in P_4\) has coordinate one in each of the six pair-linking
characters of \(P_4^{\rm ab}=\mathbb Z^6\). Choose one such
character \(\ell:P_4\to\mathbb Z\), so \(\ell(z_4)=1\).
Centrality gives the literal decomposition
\[
 \begin{gathered}
 P_4=(\ker\ell)\times\langle z_4\rangle,\qquad
 g=(gz_4^{-\ell(g)})z_4^{\ell(g)},\\
 \widehat P_4=\widehat{\ker\ell}\times
                   \widehat{\mathbb Z}_{z_4},\qquad
 \widehat{P_4/\langle z_4\rangle}=\widehat{\ker\ell}.
 \end{gathered}
 \tag{A.135.21}
\]
Completion preserves this finite direct product because any finite
quotient factors through the product of its two finite factor
images, giving cofinal product kernels. The actual action is
\(-1\) on every labelled pair-linking coordinate and inverts
\(z_4\). Thus \(\ker\ell\) is invariant and (A.135.21) is
equivariant. The pure permutation kernel of the discrete source is
exactly \(P_4/\langle z_4\rangle\); its finite index gives its
full topology in the source. Its completion is the actual \(P_T\).
Equation (A.135.20) now proves \(P_T^{\theta_T}=1\).
This step uses no general fixed-lifting principle for quotients.

\emph{5. Complete source torsion input for the whole fixed subgroup.}\par
We supply the needed source torsion argument explicitly. Normalize
the two distinguished sphere points to zero and infinity and the
first of the other three to one. The pure source is the fundamental
group of
\[
 \mathcal Z=\{(a_0,b_0)\in(\mathbb C-\{0,1\})^2:
                                      a_0\ne b_0\}.
 \tag{A.135.22}
\]
The discrete source \(\mathcal T\) is the orbifold lift group of
the action of \(S_3\) permuting \((1,a_0,b_0)\) and renormalizing
its first entry. Forgetting \(b_0\) gives free rank-two base and
free rank-three fiber. The extension splits because the base is
free. Both spaces are aspherical, so the universal cover of
\(\mathcal Z\) is contractible and
\(P_0=\pi_1\mathcal Z=F_3\rtimes F_2\).

Characteristic finite-index refinements in the finitely generated
free fiber give full induced topology and the exact completed
split sequence. They also give residual finiteness, and
characteristic refinements in \(P_0\) give residual finiteness of
its finite extension \(\mathcal T\). The discrete and profinite
Lyndon--Hochschild--Serre spectral sequences have isomorphic
finite-module terms: free groups are good, their finite-module
cohomology is finite, and only the fiber rows zero and one occur.
This proves goodness of \(P_0\); comparison for the finite
\(S_3\) extension proves goodness of \(\mathcal T\).
The free-by-free finite graph-of-spaces resolution has finite rank
in each degree. Passing through filtered unions gives the discrete
continuous-module comparison needed in A.124's prime-order input.
In particular \(P_0\) has cohomological dimension at most two,
its full completion is torsion-free, and \(\mathcal T\) has the
required finite virtual \(p\)-cohomological type.
Thus the same fully checked prime-order theorem used in A.124,
\citep[Theorem~A(i)--(iii)]{BoggiZalesskii2025}, applies to this
particular source and its full completion.

A three-cycle acts on (A.135.22) by
\((a_0,b_0)\mapsto(b_0/a_0,1/a_0)\). Its fixed points are
exactly \((\zeta,\zeta^2)\) and \((\zeta^2,\zeta)\), where
\(\zeta\) is a primitive cube root. A transposition has no fixed
point: swapping the last two points would collide them; swapping
the first two would give \((1/a_0,b_0/a_0)=(a_0,b_0)\), forcing
\(a_0=1\) since \(b_0\ne0\). The third case is conjugate.

For completeness, choose an equivariant locally finite
triangulation with these isolated fixed points as vertices and
subdivide barycentrically. Lift it to the contractible
four-dimensional universal cover. A prime-order subgroup of the
orbifold lift group injects into \(S_3\). Its positive-dimensional
cells have trivial stabilizers, since their images downstairs
would otherwise supply positive-dimensional fixed sets. Let
\(I\) be its fixed vertices. Apply a finite-rank free resolution
of \(C_p\) to cellular cochains over \(\mathbb F_p\).
Contractibility identifies one total cohomology with
\(H^*(C_p,\mathbb F_p)\). In positive group degree, regular
cell orbits are coinduced and contribute nothing; only the
fixed vertices contribute \(\prod_I\mathbb F_p\) in horizontal
degree zero. Products are exact and the resolution has finite
rank in each degree. Above total degree four no differential
can meet the degree-zero row, so
\[
           \prod_{i\in I}\mathbb F_p
                    \cong H^n(C_p,\mathbb F_p)=\mathbb F_p
                    \qquad(n>4).
 \tag{A.135.23}
\]
Thus \(I\) consists of exactly one vertex. There are no
involutions because transpositions have no fixed point
downstairs. Every order-three subgroup has exactly one fixed
vertex, and its normalizer preserves that vertex, whose
stabilizer is \(C_3\). Both centralizer and normalizer are
therefore that \(C_3\). The two downstairs fixed points are
exchanged by a transposition and deck transformations are
transitive over each point. There is one conjugacy class of
nontrivial finite subgroups and two classes of their nonidentity
elements. The prime-order comparison transfers the absence of
involutions and this order-three classification to \(T\).
There are no further finite subgroups: they inject into \(S_3\)
because the completed pure kernel is torsion-free, and the
absence of involutions leaves only \(C_3\).

Abelianizing (A.135.10) gives
\[
 T^{\rm ab}
 =\frac{\widehat{\mathbb Z}t\oplus\widehat{\mathbb Z}s_b}
        {\overline{\langle3t+6s_b\rangle}}
 \cong\widehat{\mathbb Z}\oplus C_3,
 \qquad [R]=t+2s_b\ne0.
 \tag{A.135.24}
\]
The two nonidentity torsion classes are represented by
\(R,R^{-1}\); their nonzero \(C_3\) classes distinguish them.
The action \(\theta_T\) is minus one on (A.135.24).
Its fixed subgroup injects into \(S_3\) because its intersection
with \(P_T\) is one. Thus it is finite. If nontrivial, it would
contain an element conjugate to \(R\) or \(R^{-1}\), whose
nonzero order-three abelian class cannot equal its negative.
This proves \(T^{\theta_T}=1\). We have used abelianization only
after the complete torsion classification.

\emph{6. The original gauge and the two excluded cubics.}\par
Equation (A.135.5), with \(r_S\in P_T'\), carries both fixed
subgroups to their \(\beta_T\) counterparts. This proves
(A.135.7). It also shows that \(g\mapsto g\theta_T(g)^{-1}\)
is injective on \(T\): equality for \(g,h\) means that
\(h^{-1}g\) is fixed. This is uniqueness of an existing specified
norm solution, not existence for every cocycle.

Only at this point use A.134's full embedding of the source.
The two equivariance identities in (A.135.6) identify the fixed
subgroups in the closed image with those in the source, proving
(A.135.8). A.133 proves that \(\Pi_0\) is fixed by \(\Theta_0\)
and \(\Pi_u^*\) is fixed by the ORIGINAL \(\Theta\). Both
are nontrivial of exact order three, with their specified
Y-direction markings. They cannot belong to the corresponding
trivial fixed subgroups in (A.135.8), proving (A.135.9).

No conclusion here makes a cube fixed merely because it is sent
to a conjugate of its inverse. In particular the positive
annular word of A.134 still requires its actual cube relation;
the proved negative Y-direction cubics cannot supply it by an
unmarked source identification. The full hexagon, original
boundary, and CM value remain open.
\end{proof}

\paragraph{Source and audit record.}
The complete source-gauge calculation is retained in
\path{work_A/actual_annular_source_gauge_identity37.md}.
The independent full proof of the fixed-group transfer is in
\path{work_DE/actual_annular_negative_fixed_group37.md}.
Both entire proofs and the full current C.28 were independently
read before this proposal. The source and original-gauge comparison
notes retain their earlier, correctly qualified image statements;
the present literal source identity strengthens them without
altering their bytes. No private build, finite target computation,
or edit to the frozen checkpoint37 paper was used.
